\section{Resultados}\label{sec:resultados}

\subsection{Resultados EDA}

\subsubsection{Distribución de valores faltantes}

\begin{figure}[H]
  \centering
  \includegraphics[width=\textwidth]{support2_missing_values}
  \caption{Distribución de valores faltantes en el dataset SUPPORT2.}
  \label{fig:support2-missing-values}
\end{figure}

\subsubsection{Distribución de nulos según mortalidad hospitalaria (hospdead)}

\begin{table}[H]
  \centering
  \caption[Nulos por variable según hospdead]{Distribución de nulos por variable según desenlace hospitalario (hospdead).}
  \label{tab:nulos-hospdead}
  \small
  \begin{tabular}{lrrrrr}
    \toprule
    \textbf{Variable} & \textbf{Nulos\_Vivos} & \textbf{Nulos\_Muertos} & \textbf{Porc\_Vivos} & \textbf{Porc\_Muertos} & \textbf{Diff} \\
    \midrule
    meanbp   &    1 &    0 &  0.01 &  0.00 &  0.01 \\
    hrt      &    1 &    0 &  0.01 &  0.00 &  0.01 \\
    resp     &    1 &    0 &  0.01 &  0.00 &  0.01 \\
    temp     &    1 &    0 &  0.01 &  0.00 &  0.01 \\
    sod      &    1 &    0 &  0.01 &  0.00 &  0.01 \\
    scoma    &    0 &    1 &  0.00 &  0.04 &  0.04 \\
    sps      &    0 &    1 &  0.00 &  0.04 &  0.04 \\
    aps      &    0 &    1 &  0.00 &  0.04 &  0.04 \\
    surv2m   &    0 &    1 &  0.00 &  0.04 &  0.04 \\
    surv6m   &    0 &    1 &  0.00 &  0.04 &  0.04 \\
    alb      & 2496 &  876 & 37.01 & 37.12 &  0.11 \\
    dnr      &   16 &   14 &  0.24 &  0.59 &  0.36 \\
    dnrday   &   16 &   14 &  0.24 &  0.59 &  0.36 \\
    crea     &   56 &   11 &  0.83 &  0.47 &  0.36 \\
    race     &   23 &   19 &  0.34 &  0.81 &  0.46 \\
    avtisst  &   75 &    7 &  1.11 &  0.30 &  0.82 \\
    charges  &  111 &   61 &  1.65 &  2.58 &  0.94 \\
    wblc     &  183 &   29 &  2.71 &  1.23 &  1.48 \\
    bun      & 3193 & 1159 & 47.34 & 49.11 &  1.77 \\
    urine    & 3569 & 1293 & 52.91 & 54.79 &  1.87 \\
    glucose  & 3298 & 1202 & 48.90 & 50.93 &  2.04 \\
    bili     & 1968 &  633 & 29.18 & 26.82 &  2.36 \\
    totcst   &  612 &  276 &  9.07 & 11.69 &  2.62 \\
    prg2m    & 1160 &  489 & 17.20 & 20.72 &  3.52 \\
    prg6m    & 1146 &  487 & 16.99 & 20.64 &  3.65 \\
    adls     & 2051 &  816 & 30.41 & 34.58 &  4.17 \\
    income   & 2087 &  895 & 30.94 & 37.92 &  6.98 \\
    totmcst  & 2432 & 1043 & 36.06 & 44.19 &  8.14 \\
    edu      & 1021 &  613 & 15.14 & 25.97 & 10.84 \\
    ph       & 1891 &  393 & 28.04 & 16.65 & 11.38 \\
    pafi     & 1927 &  398 & 28.57 & 16.86 & 11.70 \\
    sfdm2    & 1367 &   33 & 20.27 &  1.40 & 18.87 \\
    adlp     & 3476 & 2165 & 51.53 & 91.74 & 40.20 \\
    \bottomrule
  \end{tabular}
\end{table}

\subsubsection{Distribución de variables}

\paragraph{Variables continuas}

\begin{figure}[H]
  \centering
  \begin{subfigure}{0.48\textwidth}
    \centering
    \includegraphics[width=\textwidth]{dis_continua_1}
    \caption{Distribución de variables continuas (I).}
    \label{fig:dis-continual-1}
  \end{subfigure}
  \hfill
  \begin{subfigure}{0.48\textwidth}
    \centering
    \includegraphics[width=\textwidth]{dis_continua_2}
    \caption{Distribución de variables continuas (II).}
    \label{fig:dis-continual-2}
  \end{subfigure}
  \caption{Distribución de variables continuas en SUPPORT2.}
  \label{fig:dis-continual}
\end{figure}

\paragraph{Variables categóricas}

\begin{figure}[H]
  \centering
  \begin{subfigure}{0.48\textwidth}
    \centering
    \includegraphics[width=\textwidth]{dis_categorical_1}
    \caption{Distribución de variables categóricas (I).}
    \label{fig:dis-categorical-1}
  \end{subfigure}
  \hfill
  \begin{subfigure}{0.48\textwidth}
    \centering
    \includegraphics[width=\textwidth]{dis_categorical_2}
    \caption{Distribución de variables categóricas (II).}
    \label{fig:dis-categorical-2}
  \end{subfigure}
  \vspace{0.5em}
  \begin{subfigure}{0.48\textwidth}
    \centering
    \includegraphics[width=\textwidth]{dis_categorical_3}
    \caption{Distribución de variables categóricas (III).}
    \label{fig:dis-categorical-3}
  \end{subfigure}
  \hfill
  \begin{subfigure}{0.48\textwidth}
    \centering
    \includegraphics[width=\textwidth]{dis_categorical_4}
    \caption{Distribución de variables categóricas (IV).}
    \label{fig:dis-categorical-4}
  \end{subfigure}
  \caption{Distribución de variables categóricas en SUPPORT2.}
  \label{fig:dis-categorical}
\end{figure}

\subsubsection{Matriz de correlación}

\begin{figure}[H]
  \centering
  \includegraphics[width=\textwidth]{correlation_matrix}
  \caption{Matriz de correlación de variables en SUPPORT2.}
  \label{fig:correlation-matrix}
\end{figure}


\subsection{Resultados de modelos supervisados}

Los resultados del benchmarking se presentan de forma sistemática según la
estrategia de preprocesamiento (V1: Bio Data; V2: Tech Data) y el esquema de
entrenamiento utilizado (validación cruzada o entrenamiento simple sin CV). En
la nueva arquitectura experimental, los modelos se entrenaron sobre el conjunto
de entrenamiento, se compararon sobre el conjunto de validación y únicamente el
mejor modelo de cada escenario fue evaluado sobre el conjunto de prueba final.
En cada escenario se reportan las métricas AUC, sensibilidad, especificidad y
accuracy, junto con las curvas ROC comparativas de todos los modelos evaluados.

\subsubsection{V1 -- Bio Data con validación cruzada}

La Tabla~\ref{tab:v1-cv-base} resume el rendimiento de los modelos entrenados
sobre el dataset clínico-estadístico mediante validación cruzada interna y
selección posterior en el conjunto de validación. GBM obtuvo el mayor AUC
(0.947), seguido por Random Forest (0.942) y Logistic Regression (0.932).

\begin{table}[H]
  \centering
  \caption[V1 CV validación]{Rendimiento comparativo de modelos supervisados --- V1, validación cruzada, conjunto de validación (df1\_Bio).}
  \label{tab:v1-cv-base}
  \small
  \begin{tabular}{llrrrr}
    \toprule
    \textbf{Dataset} & \textbf{Modelo} & \textbf{AUC} & \textbf{Sensibilidad} & \textbf{Especificidad} & \textbf{Accuracy} \\
    \midrule
    df1\_Bio & GBM                & 0.947 & 0.749 & 0.961 & 0.906 \\
    df1\_Bio & RandomForest       & 0.942 & 0.746 & 0.948 & 0.896 \\
    df1\_Bio & LogisticRegression & 0.932 & 0.709 & 0.946 & 0.885 \\
    df1\_Bio & SVM                & 0.930 & 0.723 & 0.936 & 0.881 \\
    df1\_Bio & DecisionTree       & 0.919 & 0.698 & 0.937 & 0.875 \\
    df1\_Bio & KNN                & 0.903 & 0.573 & 0.958 & 0.858 \\
    \bottomrule
  \end{tabular}
\end{table}

La mejor configuración correspondió a GBM, con 2000 árboles, min\_n = 10,
profundidad máxima de 2, tasa de aprendizaje de 0.0068 y loss\_reduction = 0.
Al evaluar este modelo seleccionado sobre el conjunto de prueba final, se obtuvo
un AUC de 0.939, sensibilidad de 0.775, especificidad de 0.954 y accuracy de
0.906 (Tabla~\ref{tab:v1-cv-test}).

\begin{table}[H]
  \centering
  \caption[V1 CV prueba]{Evaluación final del mejor modelo --- V1, validación cruzada, conjunto de prueba.}
  \label{tab:v1-cv-test}
  \small
  \begin{tabular}{llrrrr}
    \toprule
    \textbf{Modelo} & \textbf{Dataset} & \textbf{AUC} & \textbf{Sensibilidad} & \textbf{Especificidad} & \textbf{Accuracy} \\
    \midrule
    GBM & df1\_Bio & 0.939 & 0.775 & 0.954 & 0.906 \\
    \bottomrule
  \end{tabular}
\end{table}

\begin{figure}[H]
  \centering
  \includegraphics[width=0.85\textwidth]{bio_roc}
  \caption{Curvas ROC comparativas --- V1, validación cruzada, conjunto de validación.}
  \label{fig:roc-v1-cv}
\end{figure}

\subsubsection{V1 -- Bio Data sin validación cruzada}

La Tabla~\ref{tab:v1-nocv-base} recoge los resultados obtenidos mediante
entrenamiento simple, comparando los modelos en el conjunto de validación. GBM
alcanzó nuevamente el mayor AUC (0.949), seguido por Random Forest (0.943) y
SVM (0.934). En este escenario, Decision Tree mostró un comportamiento más
estable que en experimentos previos, con un AUC de 0.913.

\begin{table}[H]
  \centering
  \caption[V1 sin CV validación]{Rendimiento comparativo de modelos supervisados --- V1, entrenamiento simple, conjunto de validación (df1\_Bio).}
  \label{tab:v1-nocv-base}
  \small
  \begin{tabular}{llrrrr}
    \toprule
    \textbf{Dataset} & \textbf{Modelo} & \textbf{AUC} & \textbf{Sensibilidad} & \textbf{Especificidad} & \textbf{Accuracy} \\
    \midrule
    df1\_Bio & GBM                & 0.949 & 0.763 & 0.960 & 0.909 \\
    df1\_Bio & RandomForest       & 0.943 & 0.751 & 0.952 & 0.900 \\
    df1\_Bio & SVM                & 0.934 & 0.720 & 0.940 & 0.883 \\
    df1\_Bio & LogisticRegression & 0.933 & 0.701 & 0.950 & 0.885 \\
    df1\_Bio & DecisionTree       & 0.913 & 0.734 & 0.944 & 0.890 \\
    df1\_Bio & KNN                & 0.910 & 0.596 & 0.958 & 0.864 \\
    \bottomrule
  \end{tabular}
\end{table}

La mejor configuración correspondió a GBM, con 644 árboles, profundidad máxima
de 5, tasa de aprendizaje de 0.0215443, loss\_reduction = 0.0077426 y min\_n =
20. En el conjunto de prueba final, este modelo alcanzó un AUC de 0.942,
sensibilidad de 0.755, especificidad de 0.948 y accuracy de 0.896
(Tabla~\ref{tab:v1-nocv-test}).

\begin{table}[H]
  \centering
  \caption[V1 sin CV prueba]{Evaluación final del mejor modelo --- V1, entrenamiento simple, conjunto de prueba.}
  \label{tab:v1-nocv-test}
  \small
  \begin{tabular}{llrrrr}
    \toprule
    \textbf{Modelo} & \textbf{Dataset} & \textbf{AUC} & \textbf{Sensibilidad} & \textbf{Especificidad} & \textbf{Accuracy} \\
    \midrule
    GBM & df1\_Bio & 0.942 & 0.755 & 0.948 & 0.896 \\
    \bottomrule
  \end{tabular}
\end{table}

\begin{figure}[H]
  \centering
  \includegraphics[width=0.85\textwidth]{bio_roc_no_cv}
  \caption{Curvas ROC comparativas --- V1, entrenamiento simple, conjunto de validación.}
  \label{fig:roc-v1-nocv}
\end{figure}

\subsubsection{V2 -- Tech Data con validación cruzada}

La Tabla~\ref{tab:v2-cv-base} presenta los resultados sobre el dataset
técnico-estadístico mediante validación cruzada interna y selección en el
conjunto de validación. Random Forest y GBM alcanzaron el AUC más alto (0.978),
seguidos por SVM (0.973) y Logistic Regression (0.972), superando claramente a
los modelos de la estrategia V1 en este escenario.

\begin{table}[H]
  \centering
  \caption[V2 CV validación]{Rendimiento comparativo de modelos supervisados --- V2, validación cruzada, conjunto de validación (df2\_Tech).}
  \label{tab:v2-cv-base}
  \small
  \begin{tabular}{llrrrr}
    \toprule
    \textbf{Dataset} & \textbf{Modelo} & \textbf{AUC} & \textbf{Sensibilidad} & \textbf{Especificidad} & \textbf{Accuracy} \\
    \midrule
    df2\_Tech & RandomForest       & 0.978 & 0.824 & 0.955 & 0.921 \\
    df2\_Tech & GBM                & 0.978 & 0.815 & 0.957 & 0.920 \\
    df2\_Tech & SVM                & 0.973 & 0.827 & 0.956 & 0.922 \\
    df2\_Tech & LogisticRegression & 0.972 & 0.824 & 0.956 & 0.921 \\
    df2\_Tech & DecisionTree       & 0.966 & 0.812 & 0.945 & 0.910 \\
    df2\_Tech & KNN                & 0.958 & 0.736 & 0.952 & 0.895 \\
    \bottomrule
  \end{tabular}
\end{table}

La mejor configuración correspondió a Random Forest, con mtry = 15 y min\_n =
2. En la evaluación final sobre el conjunto de prueba, este modelo obtuvo un AUC
de 0.974, sensibilidad de 0.849, especificidad de 0.957 y accuracy de 0.928
(Tabla~\ref{tab:v2-cv-test}).

\begin{table}[H]
  \centering
  \caption[V2 CV prueba]{Evaluación final del mejor modelo --- V2, validación cruzada, conjunto de prueba.}
  \label{tab:v2-cv-test}
  \small
  \begin{tabular}{llrrrr}
    \toprule
    \textbf{Modelo} & \textbf{Dataset} & \textbf{AUC} & \textbf{Sensibilidad} & \textbf{Especificidad} & \textbf{Accuracy} \\
    \midrule
    RandomForest & df2\_Tech & 0.974 & 0.849 & 0.957 & 0.928 \\
    \bottomrule
  \end{tabular}
\end{table}

\begin{figure}[H]
  \centering
  \includegraphics[width=0.85\textwidth]{tech_roc}
  \caption{Curvas ROC comparativas --- V2, validación cruzada, conjunto de validación.}
  \label{fig:roc-v2-cv}
\end{figure}

\subsubsection{V2 -- Tech Data sin validación cruzada}

La Tabla~\ref{tab:v2-nocv-base} muestra los resultados con entrenamiento simple
y comparación en el conjunto de validación. GBM obtuvo el mayor AUC (0.978),
seguido muy de cerca por Random Forest (0.977), SVM (0.973) y Logistic
Regression (0.973). La proximidad entre los modelos principales sugiere una alta
capacidad predictiva de la representación técnico-estadística del dataset.

\begin{table}[H]
  \centering
  \caption[V2 sin CV validación]{Rendimiento comparativo de modelos supervisados --- V2, entrenamiento simple, conjunto de validación (df2\_Tech).}
  \label{tab:v2-nocv-base}
  \small
  \begin{tabular}{llrrrr}
    \toprule
    \textbf{Dataset} & \textbf{Modelo} & \textbf{AUC} & \textbf{Sensibilidad} & \textbf{Especificidad} & \textbf{Accuracy} \\
    \midrule
    df2\_Tech & GBM                & 0.978 & 0.821 & 0.955 & 0.920 \\
    df2\_Tech & RandomForest       & 0.977 & 0.832 & 0.957 & 0.924 \\
    df2\_Tech & SVM                & 0.973 & 0.830 & 0.957 & 0.924 \\
    df2\_Tech & LogisticRegression & 0.973 & 0.838 & 0.953 & 0.923 \\
    df2\_Tech & DecisionTree       & 0.966 & 0.790 & 0.944 & 0.904 \\
    df2\_Tech & KNN                & 0.962 & 0.741 & 0.954 & 0.898 \\
    \bottomrule
  \end{tabular}
\end{table}

La mejor configuración correspondió a GBM, con 200 árboles, profundidad máxima
de 6, tasa de aprendizaje de 0.0464159, loss\_reduction = 0.0464159 y min\_n =
12. En el conjunto de prueba final, este modelo alcanzó un AUC de 0.975,
sensibilidad de 0.841, especificidad de 0.950 y accuracy de 0.920
(Tabla~\ref{tab:v2-nocv-test}).

\begin{table}[H]
  \centering
  \caption[V2 sin CV prueba]{Evaluación final del mejor modelo --- V2, entrenamiento simple, conjunto de prueba.}
  \label{tab:v2-nocv-test}
  \small
  \begin{tabular}{llrrrr}
    \toprule
    \textbf{Modelo} & \textbf{Dataset} & \textbf{AUC} & \textbf{Sensibilidad} & \textbf{Especificidad} & \textbf{Accuracy} \\
    \midrule
    GBM & df2\_Tech & 0.975 & 0.841 & 0.950 & 0.920 \\
    \bottomrule
  \end{tabular}
\end{table}

\begin{figure}[H]
  \centering
  \includegraphics[width=0.85\textwidth]{tech_roc_no_cv}
  \caption{Curvas ROC comparativas --- V2, entrenamiento simple, conjunto de validación.}
  \label{fig:roc-v2-nocv}
\end{figure}
